\chapter{Production deployment és service lifecycle}

\noindent\textbf{Tanulási út: Gyors út: production alap.} Tanuld meg a release/activation contractot és a biztonságos deployment workflow-t.\par\medskip

\invariantblock{Biztonsági invariáns: aktiválás előtt ellenőrzés}{A teljes candidate generációt fordítsd és validáld, mielőtt aktívvá válhatna. Ha a verification vagy activation elbukik, az előző known-good generáció szolgáljon tovább.}

\diagnosticblockhu{Production diagnostic}{SEC-PROD-002}{Az alkalmazás hiányos production policyt deklarál, ezért a szigorú deployment contract nem bizonyítható.}{Deklaráld mind a négy kötelező invariantot: HTTPS required, debug disabled, HSTS required és database TLS required.}

Ez a fejezet nem ismétli meg a konfigurációs referencia és az observability részleteit. A production szempontból azt foglalja össze, hogyan lesz egy ellenőrzött release-ből futó service, hogyan történik a release-váltás, és hogyan marad kontrollált a process lifecycle. A teljes backup/restore és disaster recovery a következő fejezet önálló témája.

\section{Az alkalmazás production contractja}

A production security egy része magában a Velran programban deklarált:

\begin{lstlisting}[language=Velran]
production {
    https required;
    debug disabled;
    hsts required;
    database tls required;
}
\end{lstlisting}

Ez nem második secret/config rendszer. Követelményt deklarál; certificate/path/trusted network továbbra is server/operator config. Másképp fogalmazva: a Velran source azt mondja meg, \emph{minek kell igaznak lennie}, a trusted deployment config pedig azt, \emph{hogyan teszi ezt igazzá az adott környezet}. Startup elbukik, ha a valós topology nem teljesíti a contractot. Source reload nem változtathatja meg csendben; production security requirement módosítás restartot és explicit operator review-t igényel.

\section{Két érvényes production frissítési út}

A Velran két production-biztos workflow-t támogat. Tudatosan válassz közülük, és ne keverd a fájlrendszerre vonatkozó feltételezéseiket.

\begin{description}
  \item[Immutable release deployment] Teljes, verziózott release könyvtár készül, ellenőrzéssel és atomikus active-path váltással; restart vagy reload csak ott történik, ahol valóban szükséges. Ezt használd, ha az üzemeltetés explicit release-eket és migrációkat kezel.
  \item[Tranzakciós rolling source] A \code{reload.mode = "rolling"} aktív; a webfejlesztő közvetlenül feltöltheti, szerkesztheti, átnevezheti vagy törölheti a figyelt \code{.vrn} fájlokat. A szerver a megállapodott forrásfát fingerprinteli, live forgalomtól elkülönítve candidate-et épít és validál, majd csak teljes siker után aktivál atomikusan.
\end{description}

Mindkét módban alapinvariáns, hogy félkész forrásállapot nem válhat aktív runtime-má. Az adatbázis-migráció továbbra is explicit operátori művelet; a rolling source aktiválás nem futtat automatikusan schema-változást.

\section{Immutable release útvonal}

Ajánlott sorrend:

\begin{lstlisting}
./verify.sh
-> cargo build --locked --release
-> immutable release directory
-> velran-server --config ... --check-config
-> velran-cli migrate verify
-> approved velran-cli migrate apply
-> velran-cli migrate verify
-> atomic current symlink switch
-> controlled restart
-> /health/live
-> /health/ready
-> log / metric / audit review
\end{lstlisting}

A buildelt \code{velran-server} és \code{velran-cli} bináris, az alkalmazásforrás és a migration legyen ugyanahhoz a release-azonosítóhoz köthető. Az immutable-release workflow-ban ne egy élő \code{/srv/velran/current} könyvtár tartalmát módosítsd fájlonként; új release könyvtárat készíts, majd atomikusan válts symlinket. Ez a tiltás nem vonatkozik arra a forrásfára, amelyet szándékosan tranzakciós \code{rolling} módban üzemeltetsz: ott a candidate-generation határ adja az atomikusságot.

\section{systemd service}

A \code{velran-server} hosszú életű service process. Az \code{velran-cli} nem daemon: migration, auth-admin vagy más operátori művelet idejére indul és kilép.

\begin{lstlisting}
[Service]
User=velran
Group=velran
ExecStart=/usr/local/bin/velran-server \
  --config /usr/local/etc/velran/server.toml
Restart=on-failure
NoNewPrivileges=true
PrivateTmp=true
\end{lstlisting}

A service fusson nem privilegizált felhasználó alatt. systemd esetén a process-szintű cgroup felső korlátok -- például \code{MemoryMax}, \code{CPUQuota} és \code{TasksMax} -- tulajdonosa a systemd legyen; a Velran ne próbálja ugyanazokat a cgroup kontrollokat párhuzamosan írni. Minden erőforrás-határhoz egyértelmű authority tartozzon.

Közvetlen TLS esetén a 80/443 portokhoz csak a bind-service capabilityt (\path{CAP_NET_BIND_SERVICE}) add meg. Ha a Velran reverse proxy mögött, például \code{127.0.0.1:8080}-on figyel, erre nincs szükség.

\section{Controlled restart és readiness}

Process-szintű konfigurációmódosítás után előbb fusson \code{--check-config}, majd controlled restart. Behind-proxy módban a SIGHUP a logok újranyitása mellett tranzakciósan újraépíti a domain/application hosting állapotot, de listener, DB/Redis/auth kapcsolat, cgroup policy, logging sink és más process-szintű beállítás továbbra is restartot igényel.

Az application source változása az alább részletezett candidate-activation útvonalat használja: request handlingtől elkülönített compile/validation után atomikus domaincsere történik. Elutasított candidate esetén a korábbi generáció marad aktív.

Rolling deploymentnél csak ready instance kapjon forgalmat: indítsd el az újat, várd meg a readiness állapotot, majd vedd ki a régit. A health diagnosztika az observability fejezetben szerepel.

\section{Migration és rollback nem ugyanaz}

Az alkalmazás binárisának vagy forrásának visszaváltása egyszerű lehet, az adatbázis-séma visszaállítása viszont nem. Ezért production migrationnél preferáld az expand/contract mintát: az új release előbb backward-compatible sémát vezessen be, az inkompatibilis takarítás csak későbbi release-ben történjen.

Ne kapcsolj automatikus reverse SQL-t az alkalmazás symlink visszaállításához. Ha a migration adatot alakított át vagy törölt, a valódi recovery pont a bizonyított backup lehet.

\section{Recovery gate a deployment előtt}

Destruktív vagy schema-kompatibilitást érintő release előtt legyen ismert recovery point és friss restore-evidence. A backup tartalma, a DB/AppFs konzisztencia, az RPO/RTO, valamint a rollback és restore döntési fa a következő fejezetben szerepel. A deployment felelőssége itt az, hogy a release-váltás előtt ez a gate teljesüljön.

\section{Host-karbantartási boundary}

A Velran fail-closed módon tudja ellenőrizni az alkalmazáskonfigurációt, capabilityket és readiness állapotot, de a hostot nem adminisztrálja. Legyen explicit felelőse az operációs rendszer security frissítéseinek, az időszinkronnak, a lemez/inode kapacitásnak, a TLS-t lezáró edge tanúsítvány-megújításának, a DB/Redis karbantartásnak, valamint a firewall- és DNS-változásoknak. Kis telepítésnél elég egy rövid runbook, amely megnevezi a felelős rendszert és a renewal/patch ritmust; az nem elfogadható, ha ezek a feladatok implicit maradnak.

A host-karbantartást kezeld kontrollált change-ként: előbb ellenőrizd a kapacitást és a backup/recovery állapotot, lehetőleg rétegenként patch-elj vagy újíts meg, majd ismételd meg a readiness-, smoke-, log- és metrics ellenőrzéseket.

\section{Deploy utáni minimum ellenőrzés}

\begin{enumerate}
  \item az új process fut és ready;
  \item a publikus főoldal vagy smoke route válaszol;
  \item nincs új error spike a server logban;
  \item access logban a státuszkód-eloszlás normális;
  \item audit log írható;
  \item DB/Redis dependency metrikák nem romlottak;
  \item egy rollback döntéshez ismert a release hash és az adatbázis migration állapota.
\end{enumerate}


\section{Release-invariantok}

Production release-nél a deployment tooling részleteitől függetlenül maradjon igaz néhány alapinvariant:

\begin{itemize}
  \item a build és verification a locked dependency graphot használja;
  \item application source, migrationök, server/CLI binárisok és a jóváhagyott config verzió egy release recordhoz köthető;
  \item production secret nem kerül az immutable release fába;
  \item candidate alkalmazás csak teljes fordítás és validáció után válhat aktívvá;
  \item az aktiválás atomikus: request vagy a régi, vagy az új generációt látja, félbemaradt keveréket nem;
  \item hibás candidate esetén a korábbi known-good generáció szolgál tovább;
  \item a váltás előtt ismert a DB migration state és az application rollback kompatibilitási határa.
\end{itemize}

Operátori szabály: a candidate vagy átmegy a teljes gate-en és atomikusan aktiválódik, vagy nem aktiválódik. A natív fordítás hibája deployment hiba, soha nem fallback-jelzés.

\section{Mi reloadolható és mi igényel restartot?}

A source reload szándékosan szűkebb, mint a process reconfiguration. A kettőt külön operátori eseményként kezeld.

\begin{longtable}{L{0.28\textwidth}L{0.30\textwidth}L{0.32\textwidth}}
\toprule
Változás & Normál mechanizmus & Indok \\
\midrule
Velran application source/modulgráf & candidate compile, validation, atomic generation swap & application-local code generation \\
Közvetlen feltöltés/szerkesztés/törlés rolling módban & stabil fingerprint, candidate build, atomikus generation swap & PHP-szerű workflow release-könyvtár kezelés nélkül \\
Immutable release symlink célja & a figyelt source path candidate generationt indít & megőrzi az atomikus release layoutot \\
Logrotation támogatott topológiában & SIGHUP/log reopen & nincs application-policy változás \\
Listener/TLS/process topology & config check + controlled restart & process-level authority \\
DB/Redis/auth connection config & config check + controlled restart & connection/service lifecycle \\
Cgroup/resource authority & controlled restart & process-level isolation boundary \\
Production security requirement & operator review + restart & source reload nem gyengítheti csendben a policyt \\
\bottomrule
\end{longtable}

Ez a különbség megakadályozza, hogy a kényelmes source reload implicit process-config csatornává váljon.

\section{Candidate aktiválás: az operátori szabály}

Kezeld az aktiválást kétlépcsős határként: \emph{előbb compile/validate, utána activate}. Hiba esetén a known-good generáció szolgál tovább. Tranzakciós \code{rolling} módban a webfejlesztő közvetlenül feltöltheti, szerkesztheti, átnevezheti vagy törölheti a figyelt \code{.vrn} fájlokat; a szerver az így létrejövő forrásfát candidate-ként kezeli, és nem az aktív runtime-ot módosítja helyben. Alkalmazásfejlesztőként ennél a contractnál megállhatsz.

A forrástörlés ugyanennek a tranzakciónak a része. Ha egy modult, handlert vagy route-ot minden rá mutató hivatkozással együtt kivonnak, a következő érvényes generation egyszerűen már nem szolgálja ki azt a route-ot. Ha bármely élő forrás továbbra is a hiányzó modulra vagy szimbólumra hivatkozik, a candidate validáció elbukik és az előző generation marad aktív. A törlés tehát vagy egy teljes, érvényes generation részeként commitolódik, vagy nincs hatása az élő forgalomra.

\subsection*{A motorháztető alatt: hogyan őrzi meg a natív aktiválás a határt}

Alkalmazásfejlesztőként első olvasáskor ezt a részt átugorhatod. A source-reload supervisor a logikai entrypointot és a releváns forrásfát fájlrendszer-metaadat + tartalom-fingerprint alapján figyeli. Miután a fájlkészlet a debounce időn át stabil, a fordítás a request handlingtől elkülönítve történik. A candidate ugyanazokon a route-, resource-, cache-, authentication- és hosting ellenőrzéseken megy át, mint startupkor.

Belsőleg az elfogadott forrás verified executable IR-ré, generált safe Rusttá, majd \code{rustc} által fordított immutable \code{cdylib}-bé és új runtime-generációvá válik. Csak ezek sikere után történik az adott domain runtime atomikus cseréje. A régi runtime-ot már birtokló requestek azon a generáción fejeződnek be; a későbbiek az újat használják. Sikeres kódaktiváláskor a public-cache generationök is ugyanennél a határnál lépnek előre.

Native compile vagy activation hiba esetén nincs VM/interpreter fallback. Ez implementációs részlet, de fontos operátori következménnyel: a hiba látható marad, és a known-good generáció marad az authority.

\section{Pre-deploy evidence}

Forgalom- vagy release-váltás előtt legyen elegendő evidence ahhoz, hogy később négy kérdésre válaszolni lehessen:

\begin{enumerate}
  \item \textbf{Mi épült?} release identifier, source revision/artifact identity és hashek.
  \item \textbf{Milyen dependency és capability lett jóváhagyva?} locked Cargo graph, valamint supply-chain capability/provenance verification.
  \item \textbf{Milyen sémát várunk?} migration set és sikeres migration verification.
  \item \textbf{Hogyan állunk helyre?} ismert backup/restore readiness és az application rollback kompatibilitási határa.
\end{enumerate}

Ez az evidence szándékosan secretmentes. Provenance-ot és operációs állapotot bizonyít anélkül, hogy credentialt másolna a release recordba.

\section{Deployment anti-patternök}

Kerüld ezeket a shortcutokat:

\begin{itemize}
  \item unlocked release build, amely csendben dependency resolutiont módosíthat;
  \item fájlok közvetlen szerkesztése egy immutable aktív release könyvtárban úgy, hogy a tranzakciós \code{rolling} supervisort megkerülöd;
  \item application rollback automatikus adatbázis-rollbackként kezelése;
  \item migration implicit futtatása server startup közben;
  \item secret CLI argumentba, release manifestbe vagy source-ba írása;
  \item forgalom küldése readiness siker előtt;
  \item hibás candidate elfogadása egy másik execution engine-re való fallbackkel.
\end{itemize}

\section{Gyakorlat: próbálj el egy sikertelen aktiválást}
\paragraph{Gyakorlási mód: üzemeltetési szcenárió.} Használd az előszóban meghatározott üzemeltetési módszert.

Készíts olyan candidate release-t, amely szándékosan elbukik egy pre-activation ellenőrzésen. Bizonyítsd, hogy az aktív generáció változatlan marad, a hiba pedig látható az operátornak. Ezután próbálj el egy sikeres aktiválást és rollbacket is.

\section{Tipikus hiba: a build sikerét deployabilitynek tekintjük}
A lefordított artefakt csak egy része a release evidence-nek. Konfiguráció, migration, secret, TLS anyag, resource policy, health check és activation semantics mind legyen érvényes, mielőtt traffic függ az új generációtól.

\section{Folyamatos mintaprojekt: Secure Notes}
Fagyaszd be az \path{examples/secure-notes/checkpoints/18-production/main.vrn} snapshotot production candidate-ként. Csomagold a Secure Notes-ot production candidate-ként. Rögzítsd a pontos verification parancsokat, migration állapotot, config checket, activation evidence-et és rollback pontot. A cél ismételhető release-eljárás, nem egyszeri sikeres telepítés.
